\section{Gyroscope (LSM9DS1)}

\subsection{Overview}
The LSM9DS1 gyroscope is a three-axis angular velocity sensor capable of measuring rotation in degrees per second (dps). Integrated into the Arduino Nano 33 BLE Sense , the gyroscope is essential for motion tracking, stabilization, and gesture-based applications \cite{Bosch:2024b, STM:2015}.

\subsection{Technical Specifications}
The gyroscope offers robust performance with the following configurable features \cite{STM:2015}:
\begin{itemize}
    \item \textbf{Full-scale range:} $\pm245$, $\pm500$, $\pm2000$ dps.
    \item \textbf{Sensitivity:} 8.75 mdps/LSB at $\pm245$ dps.
    \item \textbf{Output data rate (ODR):} 14.9 Hz to 952 Hz.
    \item \textbf{Power consumption:} 4.3 mA in high-performance mode.
    \item \textbf{Noise density:} 0.008 dps/$\sqrt{\text{Hz}}$.
    \item \textbf{Interface:} I\textsuperscript{2}C and SPI compatible.
    \item \textbf{Temperature range:} $-40^{\circ}$C to $+85^{\circ}$C.
\end{itemize}

\subsection{Calibration Procedure}
To ensure accurate readings, the gyroscope requires bias calibration:
\begin{enumerate}
    \item Place the sensor on a stable, vibration-free surface.
    \item Record angular velocity readings for each axis while stationary.
    \item Compute the average bias for each axis and subtract it from subsequent readings.
\end{enumerate}

\subsection{Calibration Code}
The following example calibrates the gyroscope and compensates for its bias:
\begin{lstlisting}[language=C++, caption=Gyroscope Calibration Code]
#include <Arduino_LSM9DS1.h>

float biasX = 0, biasY = 0, biasZ = 0;

void calibrateGyroscope() {
    const int numSamples = 100;
    float sumX = 0, sumY = 0, sumZ = 0;

    for (int i = 0; i < numSamples; i++) {
        if (IMU.gyroscopeAvailable()) {
            float gx, gy, gz;
            IMU.readGyroscope(gx, gy, gz);
            sumX += gx;
            sumY += gy;
            sumZ += gz;
        }
        delay(10);
    }

    biasX = sumX / numSamples;
    biasY = sumY / numSamples;
    biasZ = sumZ / numSamples;
}

void setup() {
    Serial.begin(9600);
    while (!Serial);

    if (!IMU.begin()) {
        Serial.println("Failed to initialize LSM9DS1 gyroscope!");
        while (1);
    }

    Serial.println("Calibrating gyroscope...");
    calibrateGyroscope();
    Serial.println("Calibration complete.");
}

void loop() {
    float gx, gy, gz;

    if (IMU.gyroscopeAvailable()) {
        IMU.readGyroscope(gx, gy, gz);
        gx -= biasX;
        gy -= biasY;
        gz -= biasZ;

        Serial.print("Calibrated Angular Velocity (X, Y, Z): ");
        Serial.print(gx, 2); Serial.print(" dps, ");
        Serial.print(gy, 2); Serial.print(" dps, ");
        Serial.println(gz, 2);
    }
    delay(100);
}
\end{lstlisting}

\subsection{Simple Code Example}
Here is a simple code snippet to measure and display angular velocity:
\begin{lstlisting}[language=C++, caption=Gyroscope Simple Code]
#include <Arduino_LSM9DS1.h>

void setup() {
    Serial.begin(9600);
    while (!Serial);

    if (!IMU.begin()) {
        Serial.println("Failed to initialize gyroscope!");
        while (1);
    }

    Serial.println("Gyroscope initialized.");
}

void loop() {
    float gx, gy, gz;

    if (IMU.gyroscopeAvailable()) {
        IMU.readGyroscope(gx, gy, gz);
        Serial.print("Angular Velocity (X, Y, Z): ");
        Serial.print(gx); Serial.print(" dps, ");
        Serial.print(gy); Serial.print(" dps, ");
        Serial.println(gz); Serial.print(" dps");
    }
    delay(100);
}
\end{lstlisting}

\subsection{Testing}
To validate the gyroscope's accuracy:
\begin{itemize}
    \item Ensure stationary readings are close to zero after calibration.
    \item Compare rotational outputs against a reference device.
\end{itemize}

\subsection{Further Readings}
\begin{itemize}
    \item LSM9DS1 Datasheet \cite{STM:2015}.
    \item Arduino Nano 33 BLE Sense Documentation \cite{ArduinoNano33Sense2025}.
    \item Adafruit LSM9DS1 Guide \cite{AdafruitLSM9DS1}.
\end{itemize}